\subsection{Feature dimensions and numerical spectral conventions}
\label{sec:feature_details}
The graph baseline has 13 columns. For each distance $r\in\{6,8,10,12\}$~\AA, it contains the nonself contact count and count divided by $n_p-1$. Two additional columns are local clustering coefficients at 8 and 12~\AA: among the center's contact neighbors, count present neighbor--neighbor edges using the same cutoff and divide by the number of possible pairs; neighborhoods with fewer than two neighbors receive zero. The final three columns are mean distances to the nearest 4, 8, and 12 nonself neighbors, using the available neighbors when a structure is shorter than the requested count. There are no B-factors, residue-position labels, inverse-square packing sums, or exponential kernels in this baseline.

The new spectral-output convention is numerical nullity, smallest positive eigenvalue, largest positive eigenvalue, positive mean, positive median, and positive population standard deviation, in that order. All six are computed at each declared scale or pair. A spectrum with no positive eigenvalues has zero for the five positive-spectrum summaries. A zero-dimensional cochain space has zero nullity as well. Finite positive eigenvalues near the numerical tolerance require audit because changing their classification affects the minimum and zero count simultaneously. The exact tolerances are reported by the implementation metadata and matrix validation artifacts.

Native degree-zero features contain 18 columns from radii 6/9/12~\AA. Each corrected degree-zero sheaf contributes 24 columns from alpha radii 3/4/5/6~\AA\ with support 13~\AA. The core combined dimensions are therefore 31 for graph plus native and 37 for graph plus each corrected sheaf. These dimensions are checked against feature names before fitting. For degree one, static lower-endpoint pairs $(3,3),(5,5)$, persistent pairs $(3,4),(5,6)$, and static upper-endpoint pairs $(4,4),(6,6)$ each supply 12 additional features. The corresponding combined models each have 49 features. Comparing persistent additions with upper as well as lower endpoint additions addresses the possibility that increasing scale alone accounts for a gain. This grid is not selected using prediction scores.

Eigenvalue classification uses a scale-relative threshold. For $m$ eigenvalues and machine precision $\epsilon$, the default is $\max(10^{-10},64\epsilon\max(1,m))\max_i|\lambda_i|$, with zero absolute offset. Values larger than this threshold are positive modes. Values below the negative threshold trigger an error. The Schur Gram matrix uses the same scale-relative rule for its pseudoinverse. The independent nullspace path uses a relative singular-value cutoff of $\sqrt{10^{-10}}$, reflecting the squaring of singular values in the Gram matrix. Additional dimension-dependent roundoff protection is recorded in the implementation. A zero spectrum is classified as zero without introducing an arbitrary absolute geometric scale.
